\section{部署与监控}

\subsection{Log + Replay 基础设施}

部署 offline policy 前需要准备 logging pipeline：

\begin{itemize}
    \item 同步 policy 输出与环境 reward，确保 reward trace 链接输出行为。
    \item 记录 policy 的 action entropy、Q variance，以便复现失败。
    \item 针对 high-risk state 定义 alert threshold（如 Q variance > 0.5）。
\end{itemize}

\begin{warningbox}{Deployment 最常见失误}
用“Offline policy + online fine-tuning” 代替 robust evaluation：上线后才发现 reward drift。正确做法是先用 replay simulator 复验，再逐步放开 online interaction 门槛。
\end{warningbox}

\subsection{监控仪表板}

\begin{table}[H]
\centering
\begin{tabular}{p{0.22\textwidth} p{0.26\textwidth} p{0.26\textwidth}}
\toprule
Dashboard 组件 & 指标 & 触发条件 \\
\midrule
Policy behavior & Action distribution shift & KL > 0.3 \\
Reward & Mean reward vs reference & 差异 > 5\% \\
Q stability & Q variance & > 0.4 \\
\bottomrule
\end{tabular}
\caption{部署后需追踪的关键指标}
\end{table}

\subsection{本章小结}

可靠部署依赖日志、replay 复制与监控仪表板。提前设置 alert 和 entropy 控制，可以让 offline policy 在上线后更稳。

\subsection{监控仪表可视化示例}

\begin{figure}[H]
\centering
\includegraphics[page=4,width=0.9\textwidth]{07_cs224r_offline_rl_2025.pdf}
\caption{Lecture slides 中展示的监控 dashboard 和 alert thresholds\protect\footnotemark}
\end{figure}
\footnotetext{Slides 第4页，展示典型 dashboard，中列展示 reward drift 控制。}

